% =============================================================================
% 历史与鉴定词汇表 — SUPPORT APPENDIX FOR THE READER
% =============================================================================
\cleardoublepage
\thispagestyle{empty}
\chapter*{历史与鉴定词汇表}
\addcontentsline{toc}{chapter}{历史与鉴定词汇表}
\markboth{历史与鉴定词汇表}{}

\begingroup

\noindent\textit{为浏览1915年大旱文献和巴巴塞纳收容院精神科报告（1917—1979）提供的辅助词汇。}

\vspace{0.8em}
\noindent\rule{\textwidth}{0.8pt}
\vspace{0.6em}

\noindent\textbf{1. 旱灾与塞尔唐术语（1914—1915）}

\vspace{0.4em}
\begin{itemize}[itemsep=0.2em]
  \item \textbf{阿拉加迪索（Alagadiço）收容营（1915）:} 1915年大旱期间在福塔莱萨附近建立的逃荒者收容营。本书第一部分的虚构场景以此为历史参照。
  \item \textbf{帕图（Patu）收容营（1932--1933）:} 位于塞纳多尔·蓬佩乌，与1932年大旱相关。它并非若阿金1915年路线中的收容营。
  \item \textbf{逃荒者 / 莫兰布多斯（破衣人）:} 对徒步逃离内陆、衣衫褴褛（molambos，破布）以求在极端干旱中活命的塞尔唐人的通俗贬称。
  \item \textbf{乱葬坑:} 浅层集体坑穴，1915年集中营中每天都有数百具饥荒和斑疹伤寒受害者的遗体被匿名丢入其中。
\end{itemize}

\vspace{0.8em}
\noindent\rule{\textwidth}{0.5pt}
\vspace{0.6em}

\noindent\textbf{2. 二十世纪医学与精神科术语}

\vspace{0.4em}
\begin{itemize}[itemsep=0.2em]
  \item \textbf{巴巴塞纳收容院:} 1903年建于米纳斯吉拉斯州的精神病机构。1930至1980年间，它成为对没有正式精神科诊断的边缘人群进行大规模收容的中心，导致六万多人死亡（即所谓的\textit{巴西大屠杀}事件）。
  \item \textbf{电休克治疗（ECT / 电休克）:} 向大脑施加电流以诱发癫痫发作。在收容院，这种做法持续且带有惩罚性，不使用麻醉或导电凝胶，造成烧伤以及股骨和脊柱骨折。
  \item \textbf{经眶脑叶切除术（弗里曼术式）:} 1940年代引入的外科干预，将一根金属棒（类似冰锥）从眼角内侧插入大脑，以切断额叶的连接。
  \item \textbf{恶病质:} 极端慢性营养不良和深度肌肉消耗的状态。
  \item \textbf{博氏线:} 指甲上的横向沟纹，由严重饥饿或感染发作期间指甲生长暂时中断引起。
  \item \textbf{单侧黄疸:} 仅限于单只眼睛的黄色着色，这是在对若阿金（患者1,260）的检查中记录到的一种异常法证发现。
\end{itemize}

\vspace{0.8em}
\noindent\rule{\textwidth}{0.5pt}
\vspace{0.6em}

\noindent\textbf{3. 法证与档案术语（2026）}

\vspace{0.4em}
\begin{itemize}[itemsep=0.2em]
  \item \textbf{脑白质刀:} 用于脑叶切除术的尖头金属外科器械。
  \item \textbf{筛板:} 鼻腔/眼腔顶部带孔的骨板，可通往大脑底部。
  \item \textbf{遍历小说契约:} 一种叙事装置，作品的实际结构和读者的决定模拟了日记中所描述的传染本身。
\end{itemize}

\endgroup
\clearpage
